%------------------------------------------------------------------------------%

\usepackage{apresentacao}
\usepackage{tabto}

\makeatletter
\graphicspath  {{./figs/}}
\def\input@path{{./figs/}}
\makeatother

%------------------------------------------------------------------------------%
\title   {Integrais e Séries}
\author  {\large Luis Alberto D'Afonseca}
\date    {Integrais e Séries}

%------------------------------------------------------------------------------%
\begin{document}

%------------------------------------------------------------------------------%
\begin{frame}
    \titlepage
\end{frame}

\begin{frame}{Disciplina}
  \begin{enumerate}[<+->]
    \item Integrais e Séries
    \item Apostila (na verdade um rascunho de apostila) link e QRcode
    \item Thomas ou Stewart
    \item Material online
  \end{enumerate}
\end{frame}

%------------------------------------------------------------------------------%
\section{Integração}

\begin{frame}{Integração}
  \begin{enumerate}[<+->]
    \item Cálculo da área de regiões curvas
    \item Operação inversa a derivação
    \item Integrais indefinidas
    \item Integrais definidas
    \item Integrais impróprias
    \item Cálculo de áreas e Volumes
    \item Técnicas de Integração
  \end{enumerate}
\end{frame}

%------------------------------------------------------------------------------%
\section{Séries}

\begin{frame}{Séries}
  \begin{enumerate}[<+->]
    \item Sequências numéricas \tabto{40mm} -- lista infinita de números
    \item Séries numéricas     \tabto{40mm} -- Soma dos termos de uma sequência
    \item Séries de potências  \tabto{40mm} -- ``Polinômio de grau infinito''
    \item Séries de Taylor     \tabto{40mm} -- Como representar funções com séries de potências
  \end{enumerate}
\end{frame}

%------------------------------------------------------------------------------%
\section{O Método Matemático}

%------------------------------------------------------------------------------%
\begin{frame}{O Método (Secreto) da Matemática}
  \begin{enumerate}[<+->]
    \item Não basta encontrar a solução correta, \\ queremos a solução \emph{comprovadamente correta}
    \item Uma aproximação não é uma solução
    \item Não podemos usar o método empírico
    \item Toda afirmação precisa ser \emph{comprovadamente verdadeira}
    \begin{enumerate}
      \item Axiomas -- Assumimos que é verdade
      \item Proposições, Lemas, Teoremas, Corolários -- Demonstramos que é verdade
      \item Nem sempre o \textit{``obvio''} é verdadeiro
    \end{enumerate}
    \item Cada passagem em um desenvolvimento deve \emph{obrigatoriamente}
    estar ancorada em um resultado demonstrado (Por isso eles são numerados)
  \end{enumerate}
\end{frame}

%------------------------------------------------------------------------------%
\begin{frame}
  \tableofcontents
\end{frame}

%------------------------------------------------------------------------------%
\end{document}
%------------------------------------------------------------------------------%
